\addcontentsline{toc}{chapter}{KATA PENGANTAR}

\begin{center}
  \Large\textbf{KATA PENGANTAR}
\end{center}
\vspace{2ex}

% Ubah paragraf-paragraf berikut sesuai dengan yang ingin diisi pada kata pengantar

Puji dan syukur ke hadirat Allah SWT atas segala rahmat, hidayah, dan karunia -Nya sehingga penulis dapat menyelesaikan salah satu kewajiban sebagai mahasiswa departemen teknik informatika, yaitu kerja praktik dengan judul "PENGEMBANGAN MODUL KOREKSI DAN KLASIFIKASI ARTEFAK CITRA OTAK SEBAGAI BAGIAN DARI SISTEM \textit{FAULT TOLERANCE IMAGE CAPTIONING}" di Laboratorium Manajemen Cerdas Informasi.

Penulis menyadari masih banyak kekurangan baik dalam pelaksanaan kerja praktik maupun penyusunan buku laporan ini. Namun penulis berharap buku laporan ini dapat menjadi refleksi dan evaluasi bagi penulis serta dapat menambah wawasan pembaca dan menjadi sumber referensi.

Melalui buku ini, penulis mengucapkan terima kasih kepada orang-orang yang telah membantu dan memberikan dukungan dalam penyusunan laporan kerja praktik ini, baik secara langsung maupun tidak langsung kepada:

\begin{enumerate}[nolistsep]

  \item Kedua orang tua penulis.

  \item Ibu Shintami Chusnul Hidayati, S.Kom, M.Sc., Ph.D, selaku dosen pembimbing kerja praktik. 

  \item Ibu Dr. Kelly Rossa Sungkono, S.Kom., M.Kom., selaku pembimbing lapangan kerja praktik.
  
  \item Mas Ibad dan Mbak Sabrina, yang telah membantu dan membimbing penulis dalam melaksanakan kerja praktik.

\end{enumerate}


\begin{flushright}
  \begin{tabular}[b]{c}
    % Ubah kalimat berikut sesuai dengan tempat, bulan, dan tahun penulisan
    Surabaya, Agustus 2025
    \\
    \\
    \\
    \\
    Alief Gilang Permana Putra
  \end{tabular}
\end{flushright}
